\section{Conclusiones}

Este trabajo demostró que los modelos de aprendizaje automático superan
significativamente a la regresión lineal para la predicción de costos
de seguros médicos. Las principales conclusiones son:

\begin{enumerate}
  \item El tabaquismo es el predictor dominante, multiplicando el costo
        esperado por un factor de 3 a 4 respecto a no fumadores.
  \item La ingeniería de features (interacción BMI-tabaquismo, edad$^2$)
        mejoró el R$^2$ en [X]\% respecto al modelo sin transformaciones.
  \item [Completar con resultado específico del mejor modelo]
  \item La solución web permite a los usuarios calcular primas estimadas
        de forma interactiva.
\end{enumerate}

Como trabajo futuro se propone incorporar variables adicionales como
historial médico y hábitos de ejercicio, así como explorar modelos
cuantílicos para estimar intervalos de predicción.
